\section*{Capítulo \ref{ch:b3:cytoskeleton-motility} --- Dinâmica do citoesqueleto e motilidade celular}

\begin{solution}{exo:b3:cytoskeleton-motility:1}
Actina: \qty{7}{nm}, actina globular, polar (farpada/pontiaguda), ATP;
córtex, protrusão, contração com a \omterm{def:b3:cytoskeleton-motility:motors}{miosina}. \omterm{def:b3:cytoskeleton-motility:filaments}{Microtúbulo}: \qty{25}{nm},
dímero de $\alpha\beta$-tubulina, polar (mais/menos), GTP; transporte
de longa distância, fuso, \omterm{def:b3:cytoskeleton-motility:cilia}{cílios}. \omterm{def:b3:cytoskeleton-motility:filaments}{Filamento intermediário}:
\qty{10}{nm}, proteínas em hélice enrolada (queratina, vimentina,
lâmina), apolar, sem nucleotídeo; resistência mecânica da célula e do
núcleo.
\end{solution}

\begin{solution}{exo:b3:cytoskeleton-motility:2}
A concentração de monômero em que uma extremidade não cresce nem
encolhe, $k_{\text{off}}/k_{\text{on}}$. As duas extremidades diferem
porque a subunidade hidrolisa seu ATP depois da incorporação, de modo
que as extremidades apresentam espécies químicas diferentes
(actina-ATP chegando à extremidade farpada, actina-ADP saindo da
pontiaguda) com afinidades diferentes. Com concentrações críticas
iguais haveria um único equilíbrio: sem treadmilling, sem fluxo e sem
meio de pagar por uma renovação dirigida.
\end{solution}

\begin{solution}{exo:b3:cytoskeleton-motility:3}
Periferia sobre \omterm{def:b3:cytoskeleton-motility:filaments}{microtúbulos}: \omterm{def:b3:cytoskeleton-motility:motors}{cinesina}, rumo à extremidade mais.
Centro: \omterm{def:b3:cytoskeleton-motility:motors}{dineína} citoplasmática, rumo à extremidade menos. Fibra de
estresse: \omterm{def:b3:cytoskeleton-motility:motors}{miosina} II, rumo à extremidade farpada da actina. \omterm{def:b3:cytoskeleton-motility:cilia}{Cílio}:
\omterm{def:b3:cytoskeleton-motility:cilia}{dineína axonemal}, rumo à extremidade menos do duplo vizinho (a base).
\end{solution}

\begin{solution}{exo:b3:cytoskeleton-motility:4}
Protrusão (Rac para os \omterm{def:b3:cytoskeleton-motility:migration}{lamelipódios}, Cdc42 para os filopódios), adesão
(\omterm{def:b3:cytoskeleton-motility:migration}{integrinas}; a jusante de Rac e de Rho), contração (Rho, por meio da
\omterm{def:b3:cytoskeleton-motility:motors}{miosina} II), retração da retaguarda (soltura das adesões, impulsionada
pela contração).
\end{solution}

\begin{solution}{exo:b3:cytoskeleton-motility:5}
$C_{c}^{+} = 1/5 = \qty{0.2}{\micro M}$, $C_{c}^{-} = 0.8/1 =
\qty{0.8}{\micro M}$; $C_{\text{ss}} = (1 + 0.8)/(5 + 1) = \qty{0.3}{\micro
M}$; $J = 5\times 0.3 - 1 = 0.5$ subunidade por segundo,
\qty{1.35}{nm/s}.
\end{solution}

\begin{solution}{exo:b3:cytoskeleton-motility:6}
$800/8 = 100$ ATP por segundo. O metro leva $1/(8\times 10^{-7}) =
1.25\times 10^{6}$ s, cerca de duas semanas, e $1.25\times 10^{8}$
passos e ATP. A difusão levaria \num{16000} anos: uma diferença de mil
milhões de vezes.
\end{solution}

\begin{solution}{exo:b3:cytoskeleton-motility:7}
$F_{\max} = 8.3\times 10^{-20}/3.6\times 10^{-8} = \qty{2.3}{pN}$. A
mesma energia espalhada por um passo mais longo dá menos força --- uma
alavanca mais longa. A \omterm{def:b3:cytoskeleton-motility:motors}{cinesina} (passo curto, \qty{6}{pN}) serve a
cargas pesadas; a \omterm{def:b3:cytoskeleton-motility:motors}{miosina} V (passo longo) cobre mais distância por ATP
e serve a um transporte rápido e leve.
\end{solution}

\begin{solution}{exo:b3:cytoskeleton-motility:8}
(a) O taxol congela os \omterm{def:b3:cytoskeleton-motility:filaments}{microtúbulos}: o fuso não consegue buscar,
capturar e corrigir, e a mitose para no ponto de checagem; a célula que
rasteja continua a protrair, mas perde a polaridade de longo prazo. (b)
A colchicina remove os \omterm{def:b3:cytoskeleton-motility:filaments}{microtúbulos}: sem fuso, parada mitótica; o
rastejamento continua sobre a actina, mas sem direção persistente. (c)
A citocalasina capeia as extremidades farpadas: sem protrusão, sem
rastejamento; a mitose prossegue, mas a citocinese falha, dando células
binucleadas. (d) Inibição de Rac: sem \omterm{def:b3:cytoskeleton-motility:migration}{lamelipódios}, e portanto sem
rastejamento; a divisão fica em grande parte inalterada.
\end{solution}

\begin{solution}{exo:b3:cytoskeleton-motility:9}
Tempo de residência da ordem do tempo de meia-recuperação, \qty{20}{s}
(constante de tempo $20/\ln 2 \approx \qty{29}{s}$); fração imóvel
\qty{10}{\%}. A polimerização na borda precisa suprir tanto o avanço
quanto o fluxo retrógrado: $1 + 1.5 = \qty{2.5}{\micro m/min}$.
\end{solution}

\begin{solution}{exo:b3:cytoskeleton-motility:10}
Uma extremidade em crescimento mantém um capuz de tubulina-GTP; atrás
dele a hidrólise produz tubulina-GDP tensionada. A perda do capuz ---
um evento aleatório cuja probabilidade sobe quando o crescimento
desacelera --- libera a tensão e a extremidade encolhe depressa até que
um novo capuz se forme. Logo acima da \omterm{thm:b3:cytoskeleton-motility:treadmilling}{concentração crítica}, cada
\omterm{def:b3:cytoskeleton-motility:filaments}{microtúbulo} está a cada momento ou com capuz e crescendo, ou sem capuz
e colapsando, de modo que os comprimentos divergem numa população longa
e numa que desaparece, em vez de se agruparem em torno de uma média. A
célula ganha uma exploração rápida de seu volume e a capacidade de
estabilizar seletivamente --- uma extremidade mais que alcance um
cinetócoro ou um sítio cortical é capturada e mantida, e as demais são
recicladas em minutos: busca e captura.
\end{solution}

\begin{solution}{exo:b3:cytoskeleton-motility:11}
Nessa escala a inércia é desprezível e as equações do fluido são
reversíveis no tempo: um movimento que refaz a si mesmo (um remo que
vai e volta) devolve o corpo ao ponto de partida, seja qual for a
velocidade de cada remada. Uma hélice em rotação nunca refaz a si mesma
--- seu movimento é quiral e contínuo --- e por isso produz empuxo
líquido. A \qty{100}{Hz}, com mil prótons por volta, $10^{5}$ prótons
por segundo por motor.
\end{solution}

\begin{solution}{exo:b3:cytoskeleton-motility:12}
Os \omterm{def:b3:cytoskeleton-motility:cilia}{cílios} das vias aéreas não conseguem bater, o muco e as bactérias
não são removidos, e as infecções voltam. Os flagelos dos
espermatozoides, construídos sobre o mesmo \omterm{def:b3:cytoskeleton-motility:cilia}{axonema}, são imóveis:
infertilidade. No embrião, os \omterm{def:b3:cytoskeleton-motility:cilia}{cílios} móveis do nó geram um fluxo para a
esquerda que estabelece o eixo esquerda--direita; sem fluxo o lado é
escolhido ao acaso, e por isso metade dos pacientes tem os órgãos
invertidos.
\end{solution}

\begin{solution}{pb:b3:cytoskeleton-motility:1}
\textbf{1.} $\qty{300}{\micro m^3} = 3\times 10^{-13}$ L; $\times
2\times 10^{-4}$ mol/L $= 6\times 10^{-17}$ mol $= 3.6\times 10^{7}$
moléculas; $1.8\times 10^{7}$ em filamentos.
\textbf{2.} $1.8\times 10^{7}\times \qty{2.7}{nm} = \qty{4.9}{cm}$ de
filamento; a \qty{0.25}{\micro m} cada, cerca de $2\times 10^{5}$
filamentos.
\textbf{3.} A proteína de capeamento bloqueia a maioria das
extremidades farpadas, e a profilina e a timosina-$\beta$4 ligam os
monômeros, de modo que a actina-ATP realmente livre fica perto da
\omterm{thm:b3:cytoskeleton-motility:treadmilling}{concentração crítica}; o crescimento se restringe às extremidades que a
célula descapeia.
\textbf{4.} $C_{\text{ss}} = 2.2/12.9 = \qty{0.17}{\micro M}$; $J = 11.6
\times 0.17 - 1.4 = 0.6$ subunidade por segundo; \qty{1.6}{nm/s} $=
\qty{0.1}{\micro m/min}$.
\textbf{5.} $1/0.1 = 10$ vezes.
\textbf{6.} \qty{1}{\micro m/min} $= \qty{16.7}{nm/s} = 6.2$ subunidades
por segundo por filamento; $\times 2\times 10^{5} = 1.2\times 10^{6}$
ATP por segundo, cerca de $0.1\,\%$ da renovação da célula.
\textbf{7.} $1/(8\times 10^{-7}) = 1.25\times 10^{6}$ s $\approx 14.5$
dias; $1.25\times 10^{8}$ passos e ATP.
\textbf{8.} $L^{2}/2D$: $(1)^{2}/(2\times 10^{-12}) = 5\times 10^{11}$ s
$\approx \num{16000}$ anos para um metro; $(10^{-5})^{2}/(2\times
10^{-12}) = 50$ s para \qty{10}{\micro m}. A difusão serve dentro de um
corpo celular; qualquer distância maior precisa de motores.
\textbf{9.} $5\times 10^{4}/6.02\times 10^{23} = 8.3\times 10^{-20}$ J
$= 20\,k_{B}T$.
\textbf{10.} $8.3\times 10^{-20}/8\times 10^{-9} = \qty{10}{pN}$;
eficiência $6/10 \approx 60\,\%$.
\textbf{11.} $F = 6\pi\times 0.01\times 10^{-7}\times 8\times 10^{-7} =
1.5\times 10^{-14}$ N $= \qty{0.015}{pN}$, quatrocentas vezes abaixo da
\omterm{prop:b3:cytoskeleton-motility:physics}{força de parada}: um motor basta e sobra contra o arrasto viscoso; as
cargas reais são obstáculos e amarras.
\textbf{12.} $10^{4}\times 100 = 10^{6}$ ATP por segundo, $0.1\,\%$ do
orçamento do neurônio: o transporte é barato. Mas os motores param no
instante em que o ATP cai, de modo que uma falha das mitocôndrias priva
a sinapse de tudo o que é feito no corpo celular --- a falha do
transporte axonal observada nas doenças neurodegenerativas.
\textbf{13.} \qty{10}{\micro m/min} $= \qty{167}{nm/s}$; $167/2.7 = 62$
subunidades por segundo por filamento.
\textbf{14.} $200\times 20 = 4000$ filamentos; $4000\times 62 = 2.5\times
10^{5}$ subunidades por segundo.
\textbf{15.} $2.5\times 10^{5}$ ATP por segundo, $0.025\,\%$ do
orçamento.
\textbf{16.} $\qty{3}{\micro m}/(\qty{10}{\micro m/min}) = \qty{0.3}{min}
= \qty{18}{s}$; $2.5\times 10^{5}\times 18 = 4.5\times 10^{6}$
subunidades na rede.
\textbf{17.} $1\,\text{pN}\times 2.7\,\text{nm} = \qty{2.7}{pN.nm} =
0.66\,k_{B}T$: uma flutuação dessa energia ocorre com probabilidade de
cerca de $e^{-0.66} \approx 0.5$ --- frestas de uma subunidade se abrem
o tempo todo, e a catraca é inteiramente plausível.
\textbf{18.} A própria polimerização empurra, sendo a bactéria a
``membrana'' na frente de sua própria rede; com a maquinaria Arp2/3 da
célula ela alcança a velocidade lamelipodial e mais --- até
\qty{0.5}{\micro m/s}, \qty{30}{\micro m/min}.
\textbf{19.} Em $C = K_{d}$, $\theta = 1/2$ e $\mathrm{d}\theta/
\mathrm{d}C = 1/(4C)$: uma diferença de \qty{2}{\%} em $C$ dá $\Delta\theta
= 0.005$; com \num{25000} receptores por lado, $125$ a mais ocupados na
frente.
\textbf{20.} Cada lado tem cerca de \num{12500} ocupados, flutuando em
$\sqrt{\num{12500}} \approx 110$, de modo que a diferença entre os dois
lados flutua em cerca de $160$: os $125$ se perdem no ruído a cada
instante. Os receptores se religam cerca de uma vez por segundo; fazer
a média ao longo de $N$ segundos reduz o ruído por $\sqrt{N}$ --- dez
segundos o levam a $50$ e o gradiente é lido.
\textbf{21.} Rac (com Cdc42) na frente, Rho na retaguarda. Com Rac
ativa em toda parte a célula protrai por todos os lados, espalha-se e
não vai a lugar nenhum.
\textbf{22.} $30/3 = \qty{10}{\micro m/min}$: como foi dado.
\textbf{23.} A protrusão é substituída pelo bolhamento movido a
pressão; a adesão e a contração continuam dependendo da actina e da
\omterm{def:b3:cytoskeleton-motility:motors}{miosina}, e por isso a actina continua essencial para o resto do ciclo.
\textbf{24.} A frente é reconstruída a cada \qty{18}{s}: um deslocamento
da atividade de Rac redireciona a polimerização dentro de um tempo de
vida da rede, a nova frente lidera, e as etapas lentas da retaguarda
seguem.
\textbf{25.} Treadmilling in vitro \qty{1.6}{nm/s}; cerca de $14.5$ dias
para percorrer o axônio; uns $2.5\times 10^{5}$ ATP por segundo para a
protrusão.
\end{solution}
